\writeSectionFrame{\presentationTitle}{Zustand}
\begin{frame}{Steckbrief}
	\begin{itemize}
 		\item Deutscher Name: Zustand
 		\item Englischer Name: State
 		\item Gruppe: Verhaltensmuster
 	\end{itemize}
\end{frame}

\begin{frame}{Zustand}
	\begin{block}{Zweck}
 		Ermögliche es einem Objekt, sein Verhalten zu ändern, wenn sich sein interner Zustand ändert. Es wird so aussehen, als ob das Objekt seine Klasse gewechselt hat. 
 	\end{block}
\end{frame}

\begin{frame}{Allgemeines Klassendiagramm}
    \begin{tikzpicture}[scale=0.8, transform shape]
        \umlclass[x=-3,y=0]{Kontext}{
        }{
        }
    
        \umlclass[x=2,y=0]{Zustand}{
        }{
            +bearbeiten()
        }

        \umlclass[x=-1,y=-3]{KonkreterZustand1}{
        }{
            +bearbeiten()
        }

        \umlclass[x=4,y=-3]{KonkreterZustand2}{
        }{
            +bearbeiten()
        }
        
        \umlaggreg{Kontext}{Zustand}
        
        \umlinherit{KonkreterZustand1}{Zustand}
        \umlinherit{KonkreterZustand2}{Zustand}
    \end{tikzpicture}
\end{frame}

\frameWithImageOnly{\BILDERPFAD/zustand1.jpg}

\tikzumlset{
    state width=1ex
}

\begin{frame}{Zustandsdiagramm Beispiel}
    \begin{tikzpicture}[scale=0.7, transform shape]
        \umlstateinitial[name=initial, x=1, y=5.2]  % Startpunkt des Diagramms
        \umlstatefinal[name=final, x=2, y=-1.9]      % Endpunkt des Diagramms
    
        \begin{umlstate}[name=waitingState, x=4, y=5]{WaitingState}
        \end{umlstate}
    
        \begin{umlstate}[name=alarmedState, x=12, y=5]{AlarmedState}
        \end{umlstate}
    
        \begin{umlstate}[name=deadState, x=7, y=-2]{DeadState}
        \end{umlstate}
    
        \begin{umlstate}[name=fightingState, x=14, y=-2]{FightingState}
        \end{umlstate}

        \begin{umlstate}[name=goingBackState, x=7, y=2]{GoingBackState}
        \end{umlstate}        

        \begin{umlstate}[name=victoriousState, x=2, y=0]{VictoriousState}
        \end{umlstate}  

        \begin{umlstate}[name=attackingState, x=14, y=2]{attackingState}
        \end{umlstate}  
    
        \umltrans{initial}{waitingState}
        \umltrans[arg=distance<5]{waitingState}{alarmedState}
        \umlHVtrans[arg=distance<3]{alarmedState}{attackingState}
        \umltrans[arg=distance<1]{attackingState}{fightingState}
        \umltrans[arg=distance>5]{attackingState}{goingBackState}
        \umltrans[arg=distance>5]{fightingState}{goingBackState}
        \umlHVtrans[arg=angekommen]{goingBackState}{waitingState}
        \umltrans[arg=hitpoints<0]{fightingState}{deadState}
        \umltrans[arg=hitpointsPlayer<0]{fightingState}{victoriousState}
        \umltrans{victoriousState}{final}
        \umltrans{deadState}{final}
    \end{tikzpicture}
\end{frame}

\begin{frame}{Klassendiagramm Beispiel}
    \begin{tikzpicture}[scale=0.7, transform shape]
        \umlclass[x=-3,y=0]{Enemy}{
        }{
        }
    
        \umlclass[x=2,y=0]{EnemyState}{
        }{
            +react()
        }

        \umlclass[x=-1,y=-3]{WaitingState}{
        }{
            +react()
        }

        \umlclass[x=4,y=-3]{AlarmedState}{
        }{
            +react()
        }

        \umlclass[x=7,y=-1]{AttackingState}{
        }{
            +react()
        }

        \umlclass[x=2,y=-5]{FightingState}{
        }{
            +react()
        }        
        
        
        \umlaggreg{Kontext}{EnemyState}
        
        \umlinherit{AlarmedState}{EnemyState}
        \umlinherit{WaitingState}{EnemyState}
        \umlinherit{AttackingState}{EnemyState}
        \umlinherit{FightingState}{EnemyState}
    \end{tikzpicture}
\end{frame}

\begin{frame}[fragile]{Zustand}
	\begin{exampleblock}{EnemyState}
     	  \begin{lstlisting}
public abstract class EnemyState {
	public abstract void react();
}     	      
     	  \end{lstlisting}
 	\end{exampleblock}
\end{frame}

\begin{frame}[fragile]{Konkreter Zustand 1}
	\begin{exampleblock}{WaitingState}
 		\begin{lstlisting}
public class WaitingState extends EnemyState {
	public void react() {
		int distance = 0;

		distance = player.getPosition() - enemy.getPosition();
		if (distance >= WAITING_RADIUS) {
			System.out.println("Feind wartet");
		} else {
			enemy.changeState(new AlarmedState(enemy, player));
		}
	}
} 		    
 		\end{lstlisting}
 	\end{exampleblock}
\end{frame}

\begin{frame}[fragile]{Konkreter Zustand 2}
	\begin{exampleblock}{AlarmedState}
 		\begin{lstlisting}
public class AlarmedState extends EnemyState {
	public void react() {
		int distance = 0;
		distance = player.getPosition() - enemy.getPosition();
		enemy.moveForward();
		
		if (distance > 0 && distance <= ALARMED_RADIUS) {
			System.out.println("Feind ist alarmiert");
		} else {
			enemy.changeState(new AttackingState(enemy, player));
		}
	}
} 		    
 		\end{lstlisting}
 	\end{exampleblock}
\end{frame}

\begin{frame}[fragile]{Kontext}
	\begin{exampleblock}{Enemy}
 		\begin{lstlisting}
public class Enemy {
	private EnemyState enemyState;
	private int position;
	private int health;
	
	public Enemy() {
		position = 0;
		health = 10;
	}
	
	public void changeState(EnemyState enemyState) {
		this.enemyState = enemyState;
	}
	
	public void doAction() {
		enemyState.react();
	}
}
 		\end{lstlisting}
 	\end{exampleblock}
\end{frame}

\frameWithImageOnly{\BILDERPFAD/zustand2.jpg}
\begin{frame}{Zustandsdiagramm Beispiel}
    \begin{tikzpicture}[scale=0.7, transform shape]
        \umlstateinitial[name=initial, x=1, y=5.2]  % Startpunkt des Diagramms
        \umlstatefinal[name=final, x=3, y=-0.7]      % Endpunkt des Diagramms
    
        \begin{umlstate}[name=readyState, x=4, y=5]{ReadyState}
        \end{umlstate}
    
        \begin{umlstate}[name=playState, x=9, y=5]{PlayState}
        \end{umlstate}
    
        \begin{umlstate}[name=pauseState, x=14, y=2]{PauseState}
        \end{umlstate}
    
        \begin{umlstate}[name=stoppedState, x=8, y=-1]{StoppedState}
        \end{umlstate}
    
        \umltrans{initial}{readyState}
        \umltrans[arg=onPlay]{readyState}{playState}
        \umlHVtrans[arg=onPause]{playState}{pauseState}
        \umlHVtrans[arg=onPlay]{pauseState}{playState}
        \umlVHtrans[arg=onstop]{pauseState}{stoppedState}
        \umltrans[arg=onstop]{playState}{stoppedState}
        \umltrans[arg=onplay]{stoppedState}{readyState}
        \umltrans[arg=onstop]{stoppedState}{final}
    \end{tikzpicture}
\end{frame}

\begin{frame}{Klassendiagramm Beispiel}
    \begin{tikzpicture}[scale=0.7, transform shape]
        \umlclass[x=-3,y=0]{MediaPlayer}{
        }{
        }
    
        \umlclass[x=2,y=0]{State}{
        }{
            +onNext(): String\\
            +onPlay(): String\\
            +onPrevious(): String\\
            +onStop(): String\\
        }

        \umlclass[x=-1,y=-3]{PauseState}{
        }{
        }

        \umlclass[x=4,y=-3]{PlayingState}{
        }{
        }

        \umlclass[x=7,y=-1]{StoppedState}{
        }{
        }

        \umlclass[x=2,y=-5]{ReadyState}{
        }{
        }        
        
        \umlaggreg{MediaPlayer}{State}
        
        \umlinherit{StoppedState}{State}
        \umlinherit{PlayingState}{State}
        \umlinherit{PauseState}{State}
        \umlinherit{ReadyState}{State}
    \end{tikzpicture}
\end{frame}

\begin{frame}[fragile]{Zustand}
	\begin{exampleblock}{State.java}
     	  \begin{lstlisting}
public abstract class State {
    private MediaPlayer player;

    public State(MediaPlayer player) {
        this.player = player;
    }

    public abstract String onStop();
    public abstract String onPlay();
    public abstract String onNext();
    public abstract String onPrevious();
}      
     	  \end{lstlisting}
 	\end{exampleblock}
\end{frame}

\begin{frame}[fragile]{Konkreter Zustand 1}
	\begin{exampleblock}{PlayingState.java}
 		\begin{lstlisting}
public class PlayingState extends State {
    PlayingState(MediaPlayer player) {
        super(player);
    }

    @Override
    public String onStop() {
    	MediaPlayer player = getPlayer();
        player.changeState(new StoppedState(player));
        player.setCurrentTrackAfterStop();
        return "Abspielen gestoppt";
    }    
 		\end{lstlisting}
 	\end{exampleblock}
\end{frame}

\begin{frame}[fragile]{Konkreter Zustand 1}
	\begin{exampleblock}{PlayingState.java}
 		\begin{lstlisting}
@Override
public String onPlay() {
    MediaPlayer player = getPlayer();
    player.changeState(new PauseState(player));
    return "Pausiert...";
}

@Override
public String onNext() {
    MediaPlayer player = getPlayer();
    return player.nextTrack();
} 
 		\end{lstlisting}
 	\end{exampleblock}
\end{frame}

\begin{frame}[fragile]{Konkreter Zustand 1}
	\begin{exampleblock}{PlayingState.java}
 		\begin{lstlisting}
@Override
public String onPrevious() {
    MediaPlayer player = getPlayer();
    return player.previousTrack();
}
 		\end{lstlisting}
 	\end{exampleblock}
\end{frame}

\begin{frame}[fragile]{Kontext}
	\begin{exampleblock}{Enemy.java}
 		\begin{lstlisting}
public class MediaPlayer {
    private State state;
    private boolean playing;
    private List<String> playlist;
    private int currentTrack = 0;

    public MediaPlayer() {
    	playlist = new ArrayList<>();
        this.state = new ReadyState(this);
        playing = true;
        for (int i = 1; i <= 12; i++) {
            playlist.add("Track " + i);
        }
    }
 		\end{lstlisting}
 	\end{exampleblock}
\end{frame}

\frameWithImageOnly{\BILDERPFAD/zustand3.jpg}
\begin{frame}{Zustandsdiagramm Beispiel}
    \begin{tikzpicture}[scale=0.7, transform shape]
    % Initial state (start)
    \umlstateinitial[name=initial, x=-1, y=0]

    % Define the states
    \begin{umlstate}[name=normalState, x=2, y=0]{NormalState}
    \end{umlstate}

    \begin{umlstate}[name=overdrawnState, x=12, y=-3]{OverdrawnState}
    \end{umlstate}

    \begin{umlstate}[name=lockedState, x=6, y=-6]{LockedState}
    \end{umlstate}    
    % Final state (end)
    \umlstatefinal[name=final, x=1, y=-4]

    % Transitions between the states
    \umltrans[arg=Konto eröffnen]{initial}{normalState}
    \umlHVtrans[arg=Geld abheben ohne ausreichendes Guthaben]{normalState}{overdrawnState}
    \umltrans[arg=Genug Geld einzahlen]{overdrawnState}{normalState}
    \umltrans[arg=Konto sperren]{overdrawnState}{lockedState}
    \umltrans[arg=Konto sperren]{normalState}{lockedState}
    \umltrans[arg=Konto schließen]{lockedState}{final}
\end{tikzpicture}
\end{frame}

\begin{frame}{Klassendiagramm Beispiel}
    \begin{tikzpicture}[scale=0.7, transform shape]
        \umlclass[x=-3,y=-5]{Bank}{
        }{
        }
    
        \umlclass[x=-3,y=0]{BankAccount}{
        }{
        }
    
        \umlclass[x=4,y=0]{AccountState}{
        }{
            +deposit(account: BankAccount, amount:double)\\
            +withdraw(account: BankAccount, amount:double)\\
        }

        \umlclass[x=-1,y=-3]{NormalState}{
        }{
        }

        \umlclass[x=4,y=-3]{OverdrawnState}{
        }{
        }

        \umlclass[x=9,y=-3]{LockedState}{
        }{
        }        
        
        \umluniassoc{BankAccount}{AccountState}
        \umlaggreg{Bank}{BankAccount}
        
        \umlinherit{LockedState}{AccountState}
        \umlinherit{OverdrawnState}{AccountState}
        \umlinherit{NormalState}{AccountState}
    \end{tikzpicture}
\end{frame}

\begin{frame}[fragile]{Zustand}
	\begin{exampleblock}{AccountState.java}
     	  \begin{lstlisting}
public interface AccountState {
    void deposit(BankAccount account, double amount);
    void withdraw(BankAccount account, double amount);
}
     	  \end{lstlisting}
 	\end{exampleblock}
\end{frame}

\begin{frame}[fragile]{Konkreter Zustand 1}
	\begin{exampleblock}{NormalState}
 		\begin{lstlisting}
public class NormalState implements AccountState {
    @Override
    public void deposit(BankAccount account, double amount) {
        account.setBalance(account.getBalance() + amount);
        System.out.println("..." + account.getBalance());
    }

    @Override
    public void withdraw(BankAccount account, double amount) {
        if (account.getBalance() >= amount) {
            account.setBalance(account.getBalance() - amount);
            System.out.println("...");
        } else {
            System.out.println("...");
            account.setState(new OverdrawnState());
            account.setBalance(account.getBalance() - amount);  
        }
    }
} 
 		\end{lstlisting}
 	\end{exampleblock}
\end{frame}

\begin{frame}[fragile]{Konkreter Zustand 1}
	\begin{exampleblock}{OverdrawnState.java}
 		\begin{lstlisting}
public class OverdrawnState implements AccountState {
    @Override
    public void deposit(BankAccount account, double amount) {
        account.setBalance(account.getBalance() + amount);
        System.out.println("...");

        if (account.getBalance() >= 0) {
            System.out.println("...");
            account.setState(new NormalState());
        }
    }

    @Override
    public void withdraw(BankAccount account, double amount) {
        System.out.println("...");
    }
}
 		\end{lstlisting}
 	\end{exampleblock}
\end{frame}


\begin{frame}[fragile]{Kontext}
	\begin{exampleblock}{BankAccount.java}
 		\begin{lstlisting}
public class BankAccount {
    private AccountState state;
    private double balance;
    private int accountId;

    public void deposit(double amount) {
        state.deposit(this, amount);
    }

    public void withdraw(double amount) {
        state.withdraw(this, amount);
    }

    public void lock() {
        setState(new LockedState());
    }
}
 		\end{lstlisting}
 	\end{exampleblock}
\end{frame}